%!TEX root=../../main.tex

As in the previous chapters, our analysis will be based on enumerating
the activation patterns associated with ReLU neurons.
However, unlike the previous chapters, we focus on training deep ReLU networks.
As a result, our description of activation patterns and the associated
constraint cones will necessarily be hierarchical.
Before proceeding, we briefly review our notation for deep networks as
introduced in \cref{sec:preliminaries}.

Recall that a deep ReLU network is
parameterized by \( \theta = \rbr{W^{(1)}, \ldots, \Wk} \), where
\( W^{(l)} \in \R^{d_{l-1} \times d_{l}} \) are the weights for the
\( l^{\text{th}} \) layer.
The dimension of the network inputs is given by \( d_0 = d \) and the output
dimension is \( d_{L} = c \), where \( c \) is the number of targets.
In accordance with \cref{eq:activations} and \cref{eq:predictions}, the model
activations and predictions are given by
\[
    Z^{(l+1)} = \rbr{\Zl W^{(l)}}_+
    \quad \text{and} \quad
    f_\theta(X) = Z^{(L)} W^{(L)},
\]
respectively, where we use \( Z^{(1)} = X \) for convenience.

Each layer of a neural network is associated with its own set of achievable
activation patterns.
In order to refer to sub-networks,
let \( \theta_l = \rbr{W^{(1)}, \ldots, \Wl} \) be the weights of the
network up to and including the \( l^{\text{th}} \) layer.
The partial network evaluation associated with this sub-network is
\begin{equation}\label{eq:sub-network}
    f_{\theta_{l}}(X) = \rbr{Z^{(l)} W^{(l)}}_+ = Z^{(l+1)}
\end{equation}
When \( l = 1 \), the sub-network is defined to be \( f_{\theta_{l}}(X) = X \).
For a given architecture specified by the hidden dimensions
\( \Delta = (d_1, \ldots, d_{L-1}) \), we can define the set of activation
patterns at layer \( l \) as follows:
\begin{equation}\label{eq:l-activation-patterns}
    \begin{aligned}
        \calDl_{\Delta} = \big\{
        \Dl_{j_l} =
        \diag\rbr{\mathds{1}\rbr{f_{\theta_{l-1}}(X) u \geq 0}}
        : \theta^{(l-1)} \in \Theta_\Delta^{(l-1)}, u \in \R^{d_{l}}
        \big\},
    \end{aligned}
\end{equation}
where \( \Theta^{(l-1)}_{\Delta} \) is the parameter space of the sub-network
associated with the architecture \( \Delta \).
When \( l = 1 \), \( \calD^{(1)}_{\Delta} \) is independent of \( \Delta \)
and equivalent to the standard definition of activation patterns for two-layer
ReLU networks given in \cref{eq:activation-patterns}.
For \( l > 1 \), the definition is dependent on the capacity of the network
as expressed by \( \Delta \).
However, it is straightforward to obtain a canonical (and maximal) set of
patterns by taking all the layer widths to be arbitrarily large.
\begin{restatable}{lemma}{limitingActivationSet}\label{lemma:activation-patterns}
    The set of activation patterns in the \( l^{\text{th}} \) layer has a
    finite limit, call it \( \calDl \), obtained by taking the hidden layer
    width \( d_{k} \rightarrow \infty \) for each \( k \in [l] \).
    \begin{equation}\label{eq:limiting-activation-set}
        \calDl = \lim_{\Delta \rightarrow \infty} \calDl_{\Delta}.
    \end{equation}
\end{restatable}
We write \( p_{l} = \abs{\calDl} \) for the number of unique activation
patterns in the \( l^{\text{th}} \) layer.

Similar to the two-layer setting, we use activation patterns to define
a data-dependent basis expansion on which the convex reformulation will
operate.
As the base case, let \( X^{(1)} = X \in \R^{n \times d_0} \) and then
recursively define the tensor 
\( X^{(l+1)} \in \R^{n \times p_1 \times \cdots \times p_{l} \times d_0} \)
in terms of the following slices:
\begin{equation}\label{eq:data-tensor}
    X^{(l+1)}_{j_1,\ldots, j_l}
    := \Dl_{j_l} \Xl_{j_1, \ldots, j_{l-1}} \in \R^{n \times d_0},
\end{equation}
where we use the convention that \( j_{l} \) always indexes the
\( l^{\text{th}} \) slice of \( \Xl \) regardless of the index
position.
Unrolling this recursion implies that \( \Xl \) is a \( l+1 \) order tensor
satisfying
\( \Xl_{j_1, \ldots, j_{l-1}} = D^{(l-1)}_{j_{l-1}} \cdots D^{(1)}_{j_1} X \).
Note that the order of the activation patterns is not important because
each \( D^{(l)}_{j_l} \) is diagonal and therefore commutative.
For clarity we shall use \( j_l \) as indices of activation patterns
and \( i_l \) as indices of neurons.

Our goal is to re-write the network \( f_\theta(X) \) as a linear tensor
operator on \( \Xk \).
To do so, we define a new tensor reduction operator as follows.
\begin{definition}[Tensor Reduction]
    Given two tensors
    \( \Xll \in \R^{n \times p_1 \times \cdots \times p_l \times d_0} \)
    and
    \( \Tll \in \R^{d_0 \times p_1 \times \cdots \times p_l \times d_{l+1}} \),
    the tensor reduction operator denoted by \( \odot \) is defined as
    \begin{equation}\label{eq:tensor-reduction}
        \Xll \odot \Tll = \sum_{j_{1}} \cdots \sum_{j_{l}}
        \Xll_{j_1, \ldots, j_{l}} \Tll_{j_1, \ldots, j_{l}}
    \end{equation}
\end{definition}
It should be clear from the definition that the tensor reduction is
linear and acts on tensors of conformal shapes.
We will also use \( \odot \) with right-hand arguments of order
\( l \), such as \( U^{(l+1)} \in \R^{d_0 \times p_1 \times \cdots \times p_{l}} \)
in which case it is given by the following reduction over
matrix-vector products:
\[
    \Xll \odot U^{(l+1)} = \sum_{j_{1}} \cdots \sum_{j_{l}}
    \Xll_{j_1, \ldots, j_{l}} U^{(l+1)}_{j_1, \ldots, j_{l}}.
\]
Using our tensor reduction operator, we define the tensor cones
associated with our higher-order activation activation patterns as
\begin{equation}\label{eq:tensor-cone}
    \calK^{(l)}_{j_{l}} =
    \cbr{
    U^{(l)} \in \R^{d_0 \times p_{1} \times \cdots \times p_{l-1}}
    :
    (2D^{(l)}_{j_{l}} - I)
    \sbr{\Xl \odot U^{(l)}} \geq 0
    }.
\end{equation}
These cones are a straightforward generalization of \( \calK_j \) from
the two-layer setting, as seen in Chapters~\ref{chapter:fast-optimization}
and~\ref{chapter:solution-functions}.
In fact, choosing \( l = 1 \) and observing that \( X^{(1)} = X \) and
\( U^{(l)} \in \R^{d_0} \) shows that \( \calK^{(1)}_{j_1} = \calK_j \).
We have already illustrated that \( \calD^{(1)} = \calD_X \); taken together,
these two facts imply that \cref{eq:tensor-cone} and
\cref{eq:l-activation-patterns} reduce directly to the standard two-layer setup
when \( L = 2 \).

Now we have the framework to prove our first result:
every neural network prediction function can be represented by a linear
operator on an \( L+1 \) order tensor.
Key to this representation is a recursive constraint set encoding the network
structure in tensor form.
Define \( \calG^{(1)} = \R^{d_0 \times d_1} \) and let \( \calGl \)
be defined recursively as
\begin{equation}\label{eq:rank-tensor-decomp}
    \begin{aligned}
         & \calGll
        := \bigg\{ \Tll
        \in \R^{d_{0} \times p_{1} \times \cdots \times p_{l} \times d_{l+1}}
        : \exists \, \Tl \in \calGl
        \text{ and } W^{(l+1)} \in \R^{d_l \times d_{l+1}} \text{ s.t. } \\
         & \hspace{0.5em} \Tll_{j_1, \ldots, j_{l}}
        = \sum_{i_l \in \calI^{(l)}_{j_l}} \Tl_{j_1, \ldots, j_{l-1}, i_{l}}
        \otimes W^{(l+1)}_{i_l},
        \calI^{(l)}_{j_{l}} = \cbr{ i_l : \Tl_{i_l} \in \calK^{(l)}_{j_{l}}},
        \sum_{j_l = 1}^{p_l} \abs{\calI^{(l)}_{j_l}} \leq d_{l}
        \bigg\},
    \end{aligned}
\end{equation}
where \( \otimes \) is the outer-product between vectors.
The set \( \calI^{(l)}_{j_l} \) contains the indices of neurons in the
\( l^{\text{th}} \) layer which map to activation pattern \( \Dl_{j_l} \),
while the constraint \( \sum_{j_l=1}^{p_l} \abs{\calI^{(l)}_{j_l}} \leq d_l \)
encodes the restriction that the number of neurons in the tensor representation
cannot exceed the layer width.
As a result, \( \Tll \in \calGll \) enforces a non-convex, low-rank structure
on the slices of \( \Tll \).

\begin{restatable}[Layer Elimination]{proposition}{layerElimination}\label{prop:layer-elimination}
    There exists \( \Tl \in \calGl \) such that
    the activations at layer \( l+2 \) are given by
    \begin{equation}\label{eq:layer-to-tensorize}
        Z^{(l+2)} =
        \rbr{\sum_{i_l=1}^{d_l} \rbr{\Xl \odot \Tl_{i_l}}_+
            \otimes W^{(l+1)}_{i_l}}_+,
    \end{equation}
    if and only if there exists \( \Tll \in \calGll \) satisfying
    \[
        Z^{(l+2)} = \rbr{X^{(l+1)} \odot \Tll}_+.
    \]
\end{restatable}

\begin{figure}[t]
    \centering
    \input{figures/chapter_four/tensor_path.tex}
    \caption{Illustration of layer elimination lemma for deep ReLU
        networks (\cref{prop:layer-elimination}).
        Each entry of the tensor \( \Tk \) is computed as the
        product of weights along the corresponding path through the neural
        network.}%
    \label{fig:tensor-path}
\end{figure}

Unrolling \cref{prop:layer-elimination} from \( l = 1 \) to \( L \)
implies that every neural network \( f_\theta \) has a tensor representation.

\begin{restatable}{theorem}{tensorEquivalence}\label{thm:tensor-equivalence}
    For each non-convex parameterization
    \( \theta = \rbr{W^{(1)}, \ldots, \Wk} \), there exists an equivalent
    tensor parameterization \( \Tk \in \calGk \) satisfying
    \[
        f_\theta(X) = \Xk \odot \Tk.
    \]
    Similarly, each tensor parameterization \( \Tk \in \calGk \)
    corresponds to a non-convex ReLU network.
\end{restatable}

As a corollary, we obtain an equivalent convex tensor problem
with non-convex decomposition constraints.

\begin{corollary}\label{cor:non-convex-tensor-program}
    The unregularized training problem for an \( L \)-layer ReLU model,
    \begin{equation}
        \min_{\theta = (W^{(1)}, \ldots, \Wk)} \calL\rbr{f_{\theta}(X), y},
    \end{equation}
    is equivalent to solving the following convex tensor program
    with non-convex decomposition constraints:
    \begin{equation}\label{eq:non-convex-tensor-program}
        \min_{\Tk}
        \calL\rbr{\Xk \odot \Tk, y} \quad \text{s.t.} \quad
        \Tk \in \calGk.
    \end{equation}
\end{corollary}

Problem~\ref{eq:non-convex-tensor-program} is attractive because it
linearizes the ReLU prediction function using the
expanded data basis given by \( \Xk \).
Moreover, it holds for any fully-connected architecture.
The drawbacks are that \( \calGk \) is non-convex and \( \Tk \) has
exponentially many parameters (recall that \( T^{(2)} \) alone
has \( d_0 p_1 c \) parameters, which is exponential when \( X \)
is full-rank).
In the next section, we derive an equivalent convex program at the cost
of new requirements on the minimum width of the network layers.

